\makeabstract
    {
    Phenocopying~\ensuremath{\Delta}acrZ~through Chromosomal Mutations in the~AcrB~Efflux Pump 
    }
    {
    Savita Jani\textsuperscript{1}, Garrett Cleveland\textsuperscript{1}, Meg Ouellette\textsuperscript{1},~Mona Wu Orr\textsuperscript{1}
    }
    {
    \textsuperscript{1} Amherst College, Amherst, MA
    }
    {
    Antibiotics are essential in treating bacterial infections. However, some bacteria have developed the ability to survive in the presence of antibiotics. This antibiotic resistance is a threat to individual and public health. In Escherichia coli, one such mechanism is the AcrAB-TolC multidrug efflux pump, a protein pump in the bacteria{\textquoteright}s cell membrane that pumps antibiotics out of the cell. The deletion of the small protein AcrZ, which is associated with the pump, changes susceptibility to some antibiotics. To better understand the specific role of AcrZ, its relation to AcrB, and its interaction with certain classes of antibiotics, we sought to phenocopy the impact of an acrZ deletion through point mutations corresponding to residues 563 and 676 in the amino acid sequence for AcrB.
    }
    {
    We first inserted the cat-sacB cassette, which confers chloramphenicol resistance and sucrose sensitivity, onto the E. coli chromosome using lambda red recombineering and allelic exchange. The cassette allowed us to test for proper chromosomal insertion using the cat-sacB phenotype, with the goal of later replacing cat-sacB with our mutants of interest.
    }
    {
    While both allelic exchange and lambda red recombineering yielded chloramphenicol resistant strains, neither method produced sucrose-sensitive strains, indicating improper integration of cat-sacB or sacB mutation.
    }
    {
    Further troubleshooting and sequencing of the strains are required before our mutants of interest can be inserted.
    }

    \newpage
    